\subsection{Scalarly essentially integrable functions}

If \(\mathbf{f}\) is a scalarly essentially \(\mu\) -integrable mapping of \(\mathrm{T}\) into \(\mathrm{F}\) , the mapping \(\mathbf{z}' \mapsto \int \langle \mathbf{f}(t), \mathbf{z}' \rangle d\mu(t)\) is a linear form on \(\mathrm{F}'\) , that is, an element of the algebraic dual \(\mathrm{F}'^*\) .

DEFINITION 1. --- One calls integral of f with respect to \(\mu\) , and denotes by \(\int f d\mu\) or \(\int f(t) d\mu(t)\) , the element of \(F'^{*}\) defined by

\[
\left\langle \mathbf {z} ^ {\prime}, \int \mathbf {f} d \mu \right\rangle = \int \left\langle \mathbf {z} ^ {\prime}, \mathbf {f} \right\rangle d \mu
\]

for all \(\mathbf{z}' \in \mathrm{F}'\) .

If f is continuous with compact support, it is scalarly integrable and Def. 1 coincides with the definition of the integral of f given in Ch. III, §3, No.~1. On the other hand, if F is a Banach space and f is essentially integrable (Ch. V, §1, No.~3, Def. 3), then f is scalarly essentially integrable and Def. 1 coincides with the definition of the integral of f given in Chapter V (Ch. V, §1, No.~3 and Ch. IV, §4, No.~2, Cor. 1 of Th. 1).

Example. --- Let X be a locally compact space, \(t \mapsto \lambda_{t}\) a mapping of T into the space \(\mathcal{M}(X)\) of measures on X. To say that the family \(t \mapsto \lambda_{t}\) is \(\mu\) -adequate means that it consists of positive measures and that the mapping \(t \mapsto \lambda_{t}\) is scalarly essentially \(\mu\) -integrable and \(\mu\) -measurable for the topology \(\sigma(\mathcal{M}(X), \mathcal{K}(X))\) .² Its integral with respect to \(\mu\) is the measure that was denoted \(\int \lambda_t d\mu(t)\) in Ch. V, §3, No.~1.

Remarks. --- 1) If F is finite-dimensional, then every scalarly essentially integrable mapping of T into F is essentially integrable (Ch. V, §1, No.~3). However, in the general case, a scalarly negligible function on a compact space T may even fail to be \(\mu\) -measurable (Exer. 12).

\begin{enumerate}
\def\labelenumi{\arabic{enumi})}
\setcounter{enumi}{1}
\tightlist
\item
  It is clear that the integral of \(\mathbf{f}\) depends only on the class of \(\mathbf{f}\) modulo the space of mappings of \(\mathbf{T}\) into \(\mathbf{F}\) that are scalarly locally \(\mu\) -negligible. Note that a scalarly locally negligible function \(\mathbf{g}\) is not necessarily zero almost everywhere (Exer. 12). However, this is indeed the case when there exists in \(\mathbf{F}'\) a sequence \((\mathbf{z}_n')\) that is dense for the topology \(\sigma(\mathbf{F}',\mathbf{F})\) : for, if \(\mathbf{H}_n\) is the locally negligible set of points \(t \in \mathbf{T}\) such that \(\langle \mathbf{g}(t), \mathbf{z}_n' \rangle \neq 0\) , the union \(\mathbf{H}\) of the \(\mathbf{H}_n\) is locally negligible and, for every \(t \notin \mathbf{H}\) , one has \(\langle \mathbf{g}(t), \mathbf{z}_n' \rangle = 0\) for all \(n\) , whence \(\mathbf{g}(t) = 0\) .
\end{enumerate}

Let \(u\) be a continuous linear mapping of \(F\) into a Hausdorff locally convex space \(G\) ; its transpose \(^{t}u\) is a linear mapping of \(G'\) into \(F'\) , and the (algebraic) transpose \(^{t}(^{t}u)\) is a linear mapping of \(F'^{*}\) into \(G'^{*}\) that extends \(u\) , which we shall again denote by \(u\) . With this convention:

PROPOSITION 1. --- If f is a mapping of T into F that is scalarly essentially μ-integrable, then the mapping u ∘ f is scalarly essentially μ-integrable and

\[
\int (u \circ \mathbf {f}) d \mu = u \Big (\int \mathbf {f} d \mu \Big).
\]

For every \(z^{\prime} \in G^{\prime}\) , \(\langle z^{\prime}, u \circ f \rangle = \langle {}^{t} u(z^{\prime}), f \rangle\) , whence the first assertion; the second follows from the formula

\[
\left\langle \mathbf {z} ^ {\prime}, \int (u \circ \mathbf {f}) d \mu \right\rangle = \int \left\langle \mathbf {z} ^ {\prime}, u \circ \mathbf {f} \right\rangle d \mu = \left\langle {} ^ {t} u (\mathbf {z} ^ {\prime}), \int \mathbf {f} d \mu \right\rangle = \left\langle \mathbf {z} ^ {\prime}, u \left(\int \mathbf {f} d \mu\right) \right\rangle .
\]

In particular, if f is scalarly essentially \(\mu\) -integrable, then it remains scalarly essentially \(\mu\) -integrable when the topology of F is replaced by a coarser topology.

PROPOSITION 2. --- Let f be a scalarly essentially \(\mu\) -integrable mapping of T into F. For every numerical function \(g \geqslant 0\) that is \(\mu\) -measurable and bounded, the mapping \(t \mapsto g(t)\mathbf{f}(t)\) (denoted gf or fg) of T into F is scalarly essentially \(\mu\) -integrable, f is scalarly essentially \((g \cdot \mu)\) -integrable, and

\[
\int \mathbf {f} d (g \cdot \mu) = \int \mathbf {f} g d \mu .
\]

This is an immediate consequence of the formula \(\langle z', gf\rangle = g\langle z', f\rangle\) , for all \(z' \in F'\) , and the formula \(\int h d(g \cdot \mu) = \int hg d\mu\) for every essentially \(\mu\) -integrable scalar function h (Ch. V, §5, No.~3, Th. 1).

A great many propositions about essentially integrable numerical functions may be transposed word for word into propositions about scalarly essentially integrable vector-valued functions. Among the more important, we call attention to the conditions for a function to be essentially integrable with respect to a measure defined by a density (Ch. V, §5, No.~3, Th. 1), or with respect to the image of a measure (Ch. V, §6, No.~2, Th. 1), or with respect to an induced measure (Ch. V, §7, No.~1, Th. 1), or with respect to the sum of a summable family of positive measures (Ch. V, §2, No.~2, Props. 1 and 3 and Cor. 3 of Prop. 1). These transpositions are left to the reader.

However, to obtain statements corresponding to the theorems on `double integrals' (Ch. V, §3, No.~3, Th. 1 and §8, No.~4, Th. 1 (Lebesgue--Fubini theorem)), it is necessary to strengthen the hypotheses (cf.~Exer. 1); on applying the previously cited theorems to each of the functions \(\langle z', f \rangle\) , where \(z' \in F'\) , one thus obtains the following propositions:

PROPOSITION 3. --- Let X be a locally compact space, \(t \mapsto \lambda_{t}\) a \(\mu\) -adequate \(^{3}\) family (Ch. V, §3, No.~1, Def. 1) of positive measures on X, and let \(\nu = \int \lambda_{t} d\mu(t)\) . Let f be a mapping of X into F; assume that \(1^{\circ} f\) is scalarly \(\nu\) -integrable; \(2^{\circ}\) there exists a locally \(\mu\) -negligible set \(N \subset T\) such that, for every \(t \notin N\) , f is scalarly \(\lambda_{t}\) -integrable and \(\int f d\lambda_{t} \in F\) . Then, the function \(t \mapsto \int f d\lambda_{t}\) , defined for \(t \notin N\) , is scalarly essentially \(\mu\) -integrable, and \(^{4}\)

\[
\int \mathbf {f} (x) d \nu (x) = \int d \mu (t) \int \mathbf {f} (x) d \lambda_ {t} (x).
\]

PROPOSITION 4. --- Let T and T' be two locally compact spaces, \(\mu\) (resp. \(\mu'\) ) a positive measure on T (resp. T'), \(\nu = \mu \otimes \mu'\) the product measure on X = T × T'. Let f be a mapping of X into F. Assume that: \(1^{\circ} f\) is scalarly \(\nu\) -integrable; \(2^{\circ}\) there exists a locally \(\mu\) -negligible set N ⊂ T such that, for every \(t \notin N\) , the mapping \(t' \mapsto f(t, t')\) is scalarly \(\mu'\) -integrable, and \(\int f(t, t') d\mu'(t') \in F\) . Then, the function \(t \mapsto \int f(t, t') d\mu'(t')\) , defined for \(t \notin N\) , is scalarly essentially \(\mu\) -integrable and \(^{4}\)

\[
\iint \mathbf {f} (t, t ^ {\prime}) d \mu (t) d \mu^ {\prime} (t ^ {\prime}) = \int d \mu (t) \int \mathbf {f} (t, t ^ {\prime}) d \mu^ {\prime} (t ^ {\prime}).
\]

